\documentclass[10pt,a4paper,twocolumn]{article}
\usepackage[utf8]{inputenc}
\usepackage[spanish]{babel}
\usepackage[margin=1.5cm]{geometry}
\usepackage{listings}
\usepackage{xcolor}
\usepackage{enumitem}

% Configuración de colores para código
\definecolor{codegreen}{rgb}{0,0.6,0}
\definecolor{codegray}{rgb}{0.5,0.5,0.5}
\definecolor{codepurple}{rgb}{0.58,0,0.82}
\definecolor{backcolour}{rgb}{0.95,0.95,0.92}

\lstdefinestyle{sqlstyle}{
    backgroundcolor=\color{backcolour},
    commentstyle=\color{codegreen},
    keywordstyle=\color{blue},
    numberstyle=\tiny\color{codegray},
    stringstyle=\color{codepurple},
    basicstyle=\ttfamily\scriptsize,
    breakatwhitespace=false,
    breaklines=true,
    captionpos=b,
    keepspaces=true,
    showspaces=false,
    showstringspaces=false,
    showtabs=false,
    tabsize=2
}


% Configuración de listings para código
\lstset{
    basicstyle=\ttfamily\footnotesize,
    keywordstyle=\color{blue}\bfseries,
    commentstyle=\color{green},
    stringstyle=\color{red},
    showstringspaces=false,
    breaklines=true,
    frame=single,
    numbers=left,
    numberstyle=\tiny\color{gray},
    inputencoding=utf8,
    % Mapeo de caracteres españoles para que no falle la compilación
    literate={á}{{\'a}}1 {é}{{\'e}}1 {í}{{\'\i}}1 {ó}{{\'o}}1 {ú}{{\'u}}1
             {Á}{{\'A}}1 {É}{{\'E}}1 {Í}{{\'I}}1 {Ó}{{\'O}}1 {Ú}{{\'U}}1
             {ñ}{{\~n}}1 {Ñ}{{\~N}}1 {ü}{{\"u}}1 {Ü}{{\"U}}1
}

\lstset{style=sqlstyle}

% Reducir el tamaño de los títulos de sección en el documento
% Usamos el paquete titlesec para controlar el formato de \section y \subsection
\usepackage{titlesec}
% Ajustes: cambiar \normalsize/\small según lo deseado
\titleformat{\section}{\normalfont\normalsize\bfseries}{\thesection}{0.5em}{}
\titleformat{\subsection}{\normalfont\small\bfseries}{\thesubsection}{0.5em}{}

\setlist[itemize]{noitemsep, topsep=0pt, leftmargin=*}
\setlength{\columnsep}{1cm}
\setlength{\parindent}{0pt}
\setlength{\parskip}{1pt}

\title{\textbf{Resumen SQL/DES}}
\author{}
\date{}

\begin{document}

\small

\maketitle

% Añadido: lista automática de secciones (TOC)
%\tableofcontents
\vspace{1em}
\section{DDL - Definición de Datos}

\textbf{CREATE TABLE}: Crea una nueva tabla.
\begin{lstlisting}
create table tabla(
  col1 tipo primary key,
  col2 tipo,
  col3 tipo
);
\end{lstlisting}

\textbf{Tipos}: \texttt{string}, \texttt{int}, \texttt{float}, etc.\\
\textbf{PRIMARY KEY}: Define clave primaria (simple o compuesta).

\section{DML - Manipulación de Datos}

\textbf{INSERT}: Inserta tuplas en una tabla.
\begin{lstlisting}
insert into tabla 
  values('val1','val2',123);
insert into tabla(col1,col2) 
  values('val1','val2');
\end{lstlisting}

\textbf{NULL}: Representa valores desconocidos/no aplicables.

\section{Consultas SELECT}

\textbf{SELECT básico}:
\begin{lstlisting}
select columnas from tabla
  where condicion;
\end{lstlisting}

\textbf{SELECT *}: Selecciona todas las columnas.\\
\textbf{DISTINCT}: Elimina duplicados en resultados.

\section{Operadores Lógicos}

\textbf{AND, OR, NOT}: Operadores booleanos.\\
\textbf{Comparación}: =, \textless\textgreater, \textless, \textgreater, \textless=, \textgreater=\\
\textbf{IS NULL / IS NOT NULL}: Comprueba valores nulos.\\
\textbf{IN / NOT IN}: Pertenencia a conjunto.

\section{VISTAS (CREATE VIEW)}

Las vistas son consultas almacenadas que pueden referenciarse como tablas.

\begin{lstlisting}
create view nombre_vista as
  select ... from ... where ...;
\end{lstlisting}

\section{Operaciones de Conjuntos}

\textbf{UNION}: Unión de resultados (elimina duplicados).
\begin{lstlisting}
select dni from programadores
union
select dni from analistas;
\end{lstlisting}

\textbf{UNION DISTINCT}: Unión eliminando duplicados explícitamente.

\textbf{INTERSECT}: Intersección de resultados.
\begin{lstlisting}
select dni from programadores
intersect
select dni from analistas;
\end{lstlisting}

\textbf{EXCEPT}: Diferencia de conjuntos (A - B).
\begin{lstlisting}
select dni from programadores
except
select dni from analistas;
\end{lstlisting}

\textbf{EXCEPT ALL}: Diferencia manteniendo duplicados.

\section{JOIN - Combinación de Tablas}

\textbf{JOIN implícito} (producto cartesiano con WHERE):
\begin{lstlisting}
select * from tabla1, tabla2
  where tabla1.id = tabla2.id;
\end{lstlisting}

\textbf{NATURAL JOIN}: Join por columnas con mismo nombre.
\begin{lstlisting}
select * from tabla1 
  natural join tabla2;
\end{lstlisting}

\textbf{JOIN explícito}:
\begin{lstlisting}
select * from tabla1 
  join tabla2 on tabla1.col1 = tabla2.col2;
\end{lstlisting}

\section{Funciones de Agregación}

\textbf{SUM(col)}: Suma de valores.\\
\textbf{COUNT(*)}: Cuenta tuplas.\\
\textbf{AVG(col)}: Media aritmética.\\
\textbf{MAX(col)}: Valor máximo.\\
\textbf{MIN(col)}: Valor mínimo.

\begin{lstlisting}
select dni, sum(horas) as total
from distribucion
group by dni;
\end{lstlisting}

\section{GROUP BY y HAVING}

\textbf{GROUP BY}: Agrupa filas por columna(s).
\begin{lstlisting}
select col, count(*) from tabla
  group by col;
\end{lstlisting}

\textbf{HAVING}: Filtra grupos (después de GROUP BY).
\begin{lstlisting}
select col, count(*) from tabla
  group by col
  having count(*) > 5;
\end{lstlisting}

\textit{Diferencia}: WHERE filtra antes de agrupar, HAVING después.

\section{Subconsultas}

\textbf{Subconsulta escalar}: Devuelve un único valor.
\begin{lstlisting}
select * from tabla where col =
  (select max(col) from tabla);
\end{lstlisting}

\textbf{Subconsulta en IN}:
\begin{lstlisting}
where dni in 
  (select dni from analistas);
\end{lstlisting}

\textbf{Subconsulta en FROM}: Tabla derivada.
\begin{lstlisting}
select * from 
  (select dni from prog 
   union select dni from anal);
\end{lstlisting}

\section{DIVISION}

Operador específico de DES para encontrar elementos que se relacionan con todos los de otro conjunto.

\begin{lstlisting}
select dniEmp from
(select codigoPr,dniEmp 
 from distribucion)
division
(select codigoPr from dist
 where dniEmp='5');
\end{lstlisting}

\textbf{Interpretación}: Empleados asignados a todos los proyectos del empleado '5'.

\section{Operaciones Aritméticas}

Se pueden realizar operaciones en SELECT:
\begin{lstlisting}
select codigo, horas*1.2 
  from distribucion;
\end{lstlisting}

\textbf{Operadores}: +, -, *, /

\section{Alias y RENAME}

\textbf{AS}: Define alias para columnas/tablas.
\begin{lstlisting}
select dni as identificador,
  sum(horas) as total
from distribucion;
\end{lstlisting}

\textbf{RENAME} (sintaxis DES):
\begin{lstlisting}
rename (dniEmp as dni)
  (distribucion);
\end{lstlisting}

\section{WITH - Expresiones de Tabla Comunes (CTE)}

\textbf{WITH ... AS} permite definir tablas temporales nombradas (CTE) que simplifican consultas complejas y mejoran legibilidad. Se pueden encadenar varias CTE y usar \texttt{WITH RECURSIVE} para recursión.

\textbf{Sintaxis básica}:
\begin{lstlisting}
WITH nombre_cte (columna1, ...) AS (
  select ...
)
SELECT ... FROM nombre_cte;
\end{lstlisting}

\textbf{Ejemplo 1: sumar horas por empleado y filtrar}
\begin{lstlisting}
WITH total_horas (dniEmp, total) AS (
  SELECT dniEmp, SUM(horas) AS total
  FROM distribución
  GROUP BY dniEmp
)
SELECT dniEmp, total
FROM total_horas
WHERE total > 20;
\end{lstlisting}

\textbf{Ejemplo 2: equivalente a DIVISION (empleados que hacen todos los proyectos de 'Evaristo')}
\begin{lstlisting}
WITH proyectos_eva (códigoPr) AS (
  SELECT códigoPr
  FROM distribución
  WHERE dniEmp = (SELECT dni FROM analistas WHERE nombre = 'Evaristo')
),
empleados_conjunto (dniEmp) AS (
  SELECT dniEmp
  FROM distribución
  WHERE códigoPr IN (SELECT códigoPr FROM proyectos_eva)
  GROUP BY dniEmp
  HAVING COUNT(DISTINCT códigoPr) = (SELECT COUNT(*) FROM proyectos_eva)
)
SELECT dniEmp FROM empleados_conjunto;
\end{lstlisting}

\section{Uso del TOP y ORDER BY}
\textbf{TOP n}: Limita el número de filas devueltas.
\textbf{ORDER BY}: Ordena resultados por columna(s), se puede usar ASC (ascendente) o DESC (descendente).
\begin{lstlisting}
SELECT TOP 5 dniEmp, SUM(horas) AS total
FROM distribución
GROUP BY dniEmp
ORDER BY total DESC;
\end{lstlisting}

\section{Ejemplos Prácticos}

\textbf{Empleados sin teléfono}:
\begin{lstlisting}
select dni,nombre from prog
where telefono is null
union
select dni,nombre from anal
where telefono is null;
\end{lstlisting}

\textbf{Suma de horas por empleado}:
\begin{lstlisting}
select dni, sum(horas) as total
from (select dni from prog 
      union select dni from anal)
join dist on dni = dniEmp
group by dni;
\end{lstlisting}

\textbf{Proyectos sin analistas}:
\begin{lstlisting}
select codigo from proyectos
except
select distinct codigo from
(select codigoPr as codigo,
        dniEmp from dist)
join analistas on dniEmp = dni;
\end{lstlisting}

\textbf{Combinando múltiples operaciones}:
\begin{lstlisting}
create view vista_compleja as
(select dni,nombre,codigoPr 
 from (select * from prog 
       union select * from anal)
 join dist on dni = dniEmp)
union
(select dni,nombre,codigo 
 from (select * from prog 
       union select * from anal)
 join proy on dni = dniDir);
\end{lstlisting}

\textbf{Combinando JOIN y WITH}:
\begin{lstlisting}
WITH pelea (barco, fecha, resultado) AS (
  SELECT barco, fecha, resultado
  FROM batallas join resultados
    ON nombre = batalla
)
SELECT barco
FROM pelea p1 JOIN pelea p2
  ON p1.barco = p2.barco
WHERE p1.fecha < p2.fecha
  AND p1.resultado = 'victoria'
  AND p2.resultado = 'derrota';
\end{lstlisting}

\textbf{Uso de 2 JOIN simultaneos}
\begin{lstlisting}
SELECT  PROJNO, DEPTNAME, MGRNO, LASTNAME
FROM    PROJECT JOIN DEPARTMENT
        ON PROJECT.DEPTNO = DEPARTMENT.DEPTNO
  JOIN EMPLOYEE
        ON DEPARTMENT.MGRNO = EMPLOYEE.EMPNO
WHERE   DEPARTMENT.DEPTNO = 'D21'
ORDER BY PROJNO;
\end{lstlisting}
Que equivale a:
\begin{lstlisting}
SELECT  PROJNO, DEPTNAME, MGRNO, LASTNAME
FROM    PROJECT P JOIN DEPARTMENT D
        ON P.DEPTNO = D.DEPTNO
  JOIN EMPLOYEE E
        ON D.MGRNO = E.EMPNO
WHERE   D.DEPTNO = 'D21'
ORDER BY P.PROJNO;
\end{lstlisting}
En este ultimo caso, se puede poner en el SELECT PROJNO o P.PROJNO, ya que no hay ambigüedad.
\end{document}